\documentclass[11pt,a4paper]{article}
\usepackage{curso}
\begin{document}
\portada{Evaluación --- Unidad}{7}{Quiz: cómputo por difusión y la analogía con Shor}
\noindent Preguntas guía: ¿por qué encontrar un orden factoriza? ¿Qué amplifica Shor y qué borra la difusión? ¿Dónde se esconde el costo?
\section*{Preguntas}
\begin{enumerate}
\item Si $r=\mathrm{ord}_N(a)$ es par y $x=a^{r/2}\not\equiv\pm1$, un factor es: (a) $\gcd(a,N)$ (b) $\gcd(x-1,N)$ (c) $x$ (d) $r$.
\item La QFT de Shor concentra la probabilidad en: (a) $y\approx kQ/r$ (b) $y=r$ (c) $y$ uniforme (d) $y=0$ solamente.
\item La caminata semi-perezosa $W=\tfrac12(I+P)$ tiene valores propios en: (a) $[-1,1]$ (b) $[0,1]$ (c) $[1,2]$ (d) $\{0,1\}$.
\item El valor del núcleo de calor en la identidad tiende a: (a) $0$ (b) $1/r$ (c) $1/N$ (d) $1$.
\item El lema de duplicación garantiza que $k2^t\bmod r$ cae en: (a) $\{0\}$ (b) $[r/4,3r/4]$ (c) $[0,r/4]$ (d) $\{r/2\}$.
\item El número de pasos del teorema de difusión es: (a) $O(r)$ (b) $O((\log N)^2)$ (c) $O(\sqrt N)$ (d) $O(1)$.
\item La lectura de una red RC con capacitores desajustados: (a) no cambia (b) cambia el valor de equilibrio (c) solo cambia la velocidad (d) da siempre $r$.
\item La primera colisión entre puntos uniformes en $r$ casillas aparece hacia: (a) $r$ (b) $\sqrt r$ (c) $\log r$ (d) $r^2$.
\item (Problema) Factorice $35$ con $a=2$ ($\mathrm{ord}_{35}(2)=12$).
\item (Problema) Calcule $\ell_0=4(M+1)(\log_2N+2)$ para $N=35$.
\item (Problema) ¿Cuántos bits de lectura harían falta para distinguir órdenes hasta $B=1000$?
\item (Problema) Dos relaciones dan $D_1=66$, $D_2=99$. ¿Qué múltiplo del orden certifica su mcd?
\end{enumerate}
\newpage
\section*{Clave}
1 (b) · 2 (a) · 3 (b) · 4 (b) · 5 (b) · 6 (b) · 7 (b) · 8 (b).
\begin{enumerate}\setcounter{enumi}{8}
\item $x=2^6\bmod35=29$; $\gcd(28,35)=7$; $35=5\cdot7$.
\item $M=6$; $\ell_0=4\cdot7\cdot(5{,}13+2)\approx200$.
\item $2\log_2 1000\approx20$ bits.
\item $\gcd(66,99)=33$: $r\mid33$ (el orden es 1, 3, 11 o 33).
\end{enumerate}
Aprobación 8/12; fallos en 4, 5 o 7 $\Rightarrow$ repasar 7.3, 7.4, 7.6.
\end{document}
